\PassOptionsToPackage{unicode}{hyperref}
\PassOptionsToPackage{hyphens}{url}
\documentclass[11pt,a4paper]{article}

% ── Codificação e língua ──────────────────────────────────────────────────────
\usepackage[T1]{fontenc}
\usepackage[utf8]{inputenc}
\usepackage[brazil]{babel}

% ── Tipografia ────────────────────────────────────────────────────────────────
\usepackage{lmodern}
\usepackage{microtype}

% ── Cores ─────────────────────────────────────────────────────────────────────
\usepackage{xcolor}
\definecolor{javaOrange}{HTML}{E76F00}
\definecolor{javaDeep}{HTML}{BF5700}
\definecolor{javaBlue}{HTML}{1565C0}
\definecolor{javaTeal}{HTML}{00695C}
\definecolor{javaAmber}{HTML}{B45309}
\definecolor{javaInk}{HTML}{2B2140}
\definecolor{javaCodeBg}{HTML}{FFF8F0}
\definecolor{javaRule}{HTML}{FFD7B5}
\definecolor{javaShade}{HTML}{FFF3E8}

% ── Caixas e código ───────────────────────────────────────────────────────────
\usepackage{tcolorbox}
\tcbuselibrary{skins,breakable,listings}
\usepackage{listings}
\usepackage{fvextra}

\lstdefinestyle{javastyle}{
  language=Java,
  basicstyle=\ttfamily\small,
  keywordstyle=\color{javaBlue}\bfseries,
  stringstyle=\color{javaDeep},
  commentstyle=\color{javaTeal}\itshape,
  numberstyle=\tiny\color{gray},
  breaklines=true,
  breakatwhitespace=false,
  showstringspaces=false,
  tabsize=4,
}

\newtcblisting{javacode}{
  listing only,
  listing style=javastyle,
  breakable,
  enhanced,
  colback=javaCodeBg,
  colframe=javaDeep,
  leftrule=4pt,
  rightrule=0.4pt,
  toprule=0.4pt,
  bottomrule=0.4pt,
  arc=2pt,
  fontupper=\small,
  top=4pt, bottom=4pt, left=6pt, right=6pt,
}

\newtcbox{\inlinecode}{
  on line,
  colback=javaCodeBg,
  colframe=javaRule,
  boxrule=0.5pt,
  arc=2pt,
  fontupper=\ttfamily\small\color{javaDeep},
  boxsep=1pt, left=2pt, right=2pt, top=1pt, bottom=1pt,
}

% ── Layout ────────────────────────────────────────────────────────────────────
\usepackage[margin=2.4cm, top=2.8cm, bottom=2.8cm]{geometry}

% ── Cabeçalho / rodapé ────────────────────────────────────────────────────────
\usepackage{fancyhdr}
\pagestyle{fancy}
\fancyhf{}
\renewcommand{\headrulewidth}{0.4pt}
\renewcommand{\headrule}{\hbox to\headwidth{\color{javaRule}\leaders\hrule height \headrulewidth\hfill}}
\fancyhead[L]{\small\sffamily\color{javaDeep} Java Programming MOOC}
\fancyhead[R]{\small\sffamily\color{javaDeep} Resoluções · Lição 2}
\fancyfoot[C]{\small\sffamily\color{gray}\thepage}

% ── Seções ────────────────────────────────────────────────────────────────────
\usepackage{titlesec}
\titleformat{\section}
  {\sffamily\Large\bfseries\color{javaOrange}}
  {\thesection}{0.7em}{}
  [\vspace{2pt}{\color{javaRule}\titlerule[1.2pt]}]
\titleformat{\subsection}
  {\sffamily\large\bfseries\color{javaDeep}}
  {\thesubsection}{0.6em}{}
\titleformat{\subsubsection}
  {\sffamily\normalsize\bfseries\color{javaTeal}}
  {\thesubsubsection}{0.5em}{}

% ── Listas ────────────────────────────────────────────────────────────────────
\usepackage{enumitem}
\setlist[itemize]{itemsep=2pt, topsep=4pt, parsep=0pt}
\setlist[enumerate]{itemsep=2pt, topsep=4pt, parsep=0pt}

% ── Tabelas ───────────────────────────────────────────────────────────────────
\usepackage{booktabs}
\usepackage{array}
\usepackage{colortbl}
\usepackage{longtable}

% ── Hiperlinks ────────────────────────────────────────────────────────────────
\usepackage{hyperref}
\hypersetup{
  colorlinks=true,
  linkcolor=javaBlue,
  urlcolor=javaBlue,
  citecolor=javaBlue,
  pdftitle={Programação em Java --- Resoluções --- Lição 2},
  pdfauthor={Java Programming MOOC · Universidade de Helsinque},
  pdflang={pt-BR},
}

% ── Caixa de destaque (dica/atenção) ─────────────────────────────────────────
\newtcolorbox{destaque}[1][Observação]{
  enhanced, breakable,
  colback=javaShade,
  colframe=javaOrange,
  leftrule=4pt, rightrule=0.4pt, toprule=0.4pt, bottomrule=0.4pt,
  arc=2pt,
  fonttitle=\sffamily\bfseries\color{javaOrange},
  title=#1,
  top=4pt, bottom=4pt,
}

% ── Título personalizado ──────────────────────────────────────────────────────
\makeatletter
\newcommand{\subtitle}[1]{\gdef\@subtitle{#1}}
\renewcommand{\maketitle}{%
  \begin{center}
    {\color{javaRule}\rule{\linewidth}{1pt}}\\[6pt]
    {\Huge\sffamily\bfseries\color{javaOrange} \@title}\\[6pt]
    \ifx\@subtitle\undefined\else
      {\Large\sffamily\color{javaDeep} \@subtitle}\\[4pt]
    \fi
    {\small\sffamily\color{gray} \@author}\\[2pt]
    {\small\sffamily\color{gray} \@date}\\[6pt]
    {\color{javaRule}\rule{\linewidth}{1pt}}
  \end{center}
  \vspace{1em}
}
\makeatother

% ─────────────────────────────────────────────────────────────────────────────
\title{Programação em Java — Resoluções}
\subtitle{Lição 2}
\author{Java Programming MOOC · Universidade de Helsinque — tradução para o português}
\date{2026}
% ─────────────────────────────────────────────────────────────────────────────

\begin{document}
\maketitle
\tableofcontents
\vspace{1.5em}

% ─────────────────────────────────────────────────────────────────────────────
\section{Exercício 1 --- Tabuada}

\begin{javacode}
import java.util.Scanner;

public class Tabuada {
    public static void main(String[] args) {
        Scanner scanner = new Scanner(System.in);

        System.out.print("Digite um número: ");
        int n = Integer.valueOf(scanner.nextLine());

        for (int i = 1; i <= 10; i++) {
            System.out.println(n + " x " + i + " = " + (n * i));
        }
    }
}
\end{javacode}

% ─────────────────────────────────────────────────────────────────────────────
\section{Exercício 2 --- Fatorial}

\begin{javacode}
public class Fatorial {

    public static int fatorial(int n) {
        int resultado = 1;
        for (int i = 2; i <= n; i++) {
            resultado *= i;
        }
        return resultado;
    }

    public static void main(String[] args) {
        System.out.println("fatorial(0) = " + fatorial(0));
        System.out.println("fatorial(1) = " + fatorial(1));
        System.out.println("fatorial(5) = " + fatorial(5));
        System.out.println("fatorial(10) = " + fatorial(10));
    }
}
\end{javacode}

\begin{destaque}
O laço começa em \inlinecode{i = 2} porque multiplicar por 1 não altera o resultado. Para \inlinecode{n = 0} e \inlinecode{n = 1}, o laço nunca executa e a função retorna \inlinecode{1}, valor correto do fatorial nesses dois casos.
\end{destaque}

% ─────────────────────────────────────────────────────────────────────────────
\section{Exercício 3 --- Contagem de pares e ímpares}

\begin{javacode}
import java.util.Scanner;

public class ContagemParesImpares {
    public static void main(String[] args) {
        Scanner scanner = new Scanner(System.in);
        int pares = 0;
        int impares = 0;

        while (true) {
            System.out.print("Digite um número (0 para parar): ");
            int numero = Integer.valueOf(scanner.nextLine());

            if (numero == 0) {
                break;
            }

            if (numero % 2 == 0) {
                pares++;
            } else {
                impares++;
            }
        }

        System.out.println("Pares: " + pares);
        System.out.println("Ímpares: " + impares);
    }
}
\end{javacode}

% ─────────────────────────────────────────────────────────────────────────────
\section{Exercício 4 --- Potência}

\begin{javacode}
public class Potencia {

    public static int potencia(int base, int expoente) {
        int resultado = 1;
        for (int i = 0; i < expoente; i++) {
            resultado *= base;
        }
        return resultado;
    }

    public static void main(String[] args) {
        System.out.println(potencia(2, 10));
    }
}
\end{javacode}

% ─────────────────────────────────────────────────────────────────────────────
\section{Exercício 5 --- Validação de entrada}

\begin{javacode}
import java.util.Scanner;

public class ValidacaoEntrada {

    public static int lerInteiroPositivo(Scanner scanner) {
        while (true) {
            System.out.print("Digite um número positivo: ");
            int numero = Integer.valueOf(scanner.nextLine());

            if (numero > 0) {
                return numero;
            }

            System.out.println("Erro: o número precisa ser positivo. Tente novamente.");
        }
    }

    public static void main(String[] args) {
        Scanner scanner = new Scanner(System.in);

        int a = lerInteiroPositivo(scanner);
        int b = lerInteiroPositivo(scanner);

        System.out.println("Produto: " + (a * b));
    }
}
\end{javacode}

% ─────────────────────────────────────────────────────────────────────────────
\section{Exercício 6 --- Estrelas}

\begin{javacode}
public class Estrelas {

    public static void imprimirEstrelas(int n) {
        String linha = "";
        for (int i = 0; i < n; i++) {
            linha += "*";
        }
        System.out.println(linha);
    }

    public static void main(String[] args) {
        for (int i = 1; i <= 5; i++) {
            imprimirEstrelas(i);
        }
    }
}
\end{javacode}

\begin{destaque}[Saída]
\begin{verbatim}
*
**
***
****
*****
\end{verbatim}
\end{destaque}

\end{document}
